\documentclass[11pt,a4paper]{article}
\usepackage[utf8]{inputenc}
\usepackage{amsmath, amssymb, amsthm}
\usepackage{physics}
\usepackage{geometry}
\geometry{margin=22mm}

\newtheorem{theorem}{Theorem}[section]
\newtheorem{lemma}[theorem]{Lemma}
\newtheorem{proposition}[theorem]{Proposition}
\newtheorem{corollary}[theorem]{Corollary}
\theoremstyle{definition}
\newtheorem{definition}[theorem]{Definition}
\theoremstyle{remark}
\newtheorem{remark}[theorem]{Remark}

\title{Rigorous Derivation of Emerging Dynamics in \\ Background-Independent Qubit Networks}
\author{JJY}
\date{\today}

\begin{document}

\maketitle

% -----------------------------------------------------------------------
\section{Introduction to the Global Null Space}
% -----------------------------------------------------------------------

The fundamental premise of our model is the existence of a global state $\ket{\Psi}$ in the
total Hilbert space $\mathcal{H}_{tot} = \mathcal{H}_C \otimes \mathcal{H}_S$.
For this mapping to be mathematically valid, the total Hamiltonian $\hat{H}_{tot}$ must possess
a non-trivial kernel (Null Space):
\begin{equation}\label{eq:kernel}
    \mathrm{Ker}(\hat{H}_{tot})
      = \bigl\{\ket{\Psi}\in\mathcal{H}_{tot}\mid\hat{H}_{tot}\ket{\Psi}=0\bigr\}
      \neq \emptyset.
\end{equation}

\begin{proposition}\label{prop:kernel}
Condition~\eqref{eq:kernel} holds \emph{if and only if} the intersection of the energy spectra
satisfies
\begin{equation}\label{eq:spec}
    \mathrm{Spec}(\hat{H}_C)\cap\mathrm{Spec}(-\hat{H}_S)\neq\emptyset.
\end{equation}
\end{proposition}

\begin{proof}
Let $\{E_k^C,\ket{e_k}_C\}$ and $\{E_j^S,\ket{f_j}_S\}$ be complete eigensystems of
$\hat{H}_C$ and $\hat{H}_S$, respectively.
Since $\hat{H}_{tot}=\hat{H}_C\otimes I_S + I_C\otimes\hat{H}_S$, the product basis
$\{\ket{e_k}_C\otimes\ket{f_j}_S\}_{k,j}$ diagonalises $\hat{H}_{tot}$:
\[
  \hat{H}_{tot}\bigl(\ket{e_k}_C\otimes\ket{f_j}_S\bigr)
    = (E_k^C+E_j^S)\bigl(\ket{e_k}_C\otimes\ket{f_j}_S\bigr).
\]

\noindent\textit{($\Rightarrow$)}\enspace
If $E_{k_0}^C+E_{j_0}^S=0$, then $\ket{e_{k_0}}_C\otimes\ket{f_{j_0}}_S\neq 0$
is a non-trivial element of $\mathrm{Ker}(\hat{H}_{tot})$.

\noindent\textit{($\Leftarrow$)}\enspace
Let $\ket{\Psi}=\sum_{k,j}c_{kj}\ket{e_k}_C\otimes\ket{f_j}_S\neq 0$ satisfy
$\hat{H}_{tot}\ket{\Psi}=0$.  Then
$\sum_{k,j}c_{kj}(E_k^C+E_j^S)\ket{e_k}_C\otimes\ket{f_j}_S = 0$.
By linear independence of the product basis, $c_{kj}(E_k^C+E_j^S)=0$ for all $k,j$.
Since $\ket{\Psi}\neq 0$ there exists $(k_0,j_0)$ with $c_{k_0j_0}\neq 0$,
forcing $E_{k_0}^C=-E_{j_0}^S\in\mathrm{Spec}(\hat{H}_C)\cap\mathrm{Spec}(-\hat{H}_S)$.
\end{proof}

In an $N$-qubit simulation, condition~\eqref{eq:spec} is achieved by tuning the local bias
terms to ensure $E_C^k+E_S^j=0$.

% -----------------------------------------------------------------------
\section{Clock Representation and Stone's Theorem}
% -----------------------------------------------------------------------

We define the \emph{clock basis} $\ket{t}_C$ through the unitary group generated by $\hat{H}_C$.
By Stone's theorem on one-parameter unitary groups, the strongly continuous map
$t\mapsto e^{-i\hat{H}_C t}$ has $\hat{H}_C$ as its unique self-adjoint generator:
\begin{equation}\label{eq:clock}
  \ket{t}_C = U(t)\ket{0}_C = e^{-i\hat{H}_C t}\ket{0}_C.
\end{equation}

Taking the Hilbert-space adjoint of~\eqref{eq:clock} gives $\bra{t}_C=\bra{0}_C e^{i\hat{H}_C t}$.
Differentiating with respect to $t$:
\begin{equation}\label{eq:dbradt}
  \frac{d}{dt}\bra{t}_C
    = \bra{0}_C\cdot i\hat{H}_C\,e^{i\hat{H}_C t}
    = \bra{0}_C\,e^{i\hat{H}_C t}\cdot i\hat{H}_C
    = \bra{t}_C(i\hat{H}_C),
\end{equation}
where the second equality uses $[e^{i\hat{H}_C t},\hat{H}_C]=0$; this commutator vanishes
because $e^{i\hat{H}_C t}=\sum_{n=0}^\infty(i\hat{H}_C t)^n/n!$ and every power of $\hat{H}_C$
commutes with $\hat{H}_C$.
Multiplying~\eqref{eq:dbradt} by $i$:
\begin{equation}\label{eq:idbra}
  i\frac{d}{dt}\bra{t}_C = -\bra{t}_C\hat{H}_C.
\end{equation}

% -----------------------------------------------------------------------
\section{Step-by-Step Flow of Unitary Emergence}
% -----------------------------------------------------------------------

We now derive the internal Schr\"odinger equation from the static WdW equation.

\subsection*{Step 1: Projection of the Global Constraint}
Starting from the WdW equation $\hat{H}_{tot}\ket{\Psi}=0$, we project onto the clock state
$\bra{t}_C$:
\begin{equation}\label{eq:proj}
  \bra{t}_C\bigl(\hat{H}_C\otimes I_S+I_C\otimes\hat{H}_S\bigr)\ket{\Psi}=0.
\end{equation}

\subsection*{Step 2: Separation of Operators}
Expanding~\eqref{eq:proj} by linearity:
\begin{equation}\label{eq:expand}
  \bra{t}_C\hat{H}_C\ket{\Psi}+\bra{t}_C(I_C\otimes\hat{H}_S)\ket{\Psi}=0.
\end{equation}
To show that $\hat{H}_S$ passes through the clock projection, write
$\ket{\Psi}=\sum_\alpha\ket{\phi_\alpha}_C\otimes\ket{\chi_\alpha}_S$.  Then
\[
  \bra{t}_C(I_C\otimes\hat{H}_S)\ket{\Psi}
    = \sum_\alpha\braket{t|\phi_\alpha}_C\,\hat{H}_S\ket{\chi_\alpha}_S
    = \hat{H}_S\!\Bigl(\sum_\alpha\braket{t|\phi_\alpha}_C\,\ket{\chi_\alpha}_S\Bigr)
    = \hat{H}_S\bigl(\bra{t}_C\ket{\Psi}\bigr),
\]
where the second equality uses the $\mathbb{C}$-linearity of $\hat{H}_S$ on $\mathcal{H}_S$
(the scalar coefficients $\braket{t|\phi_\alpha}_C\in\mathbb{C}$ factor out).
Equation~\eqref{eq:expand} becomes:
\begin{equation}\label{eq:sep}
  \bra{t}_C\hat{H}_C\ket{\Psi}+\hat{H}_S\bigl(\bra{t}_C\ket{\Psi}\bigr)=0.
\end{equation}

\subsection*{Step 3: Identification of the Relative State}
Define the \emph{relative (conditional) state} $\ket{\psi(t)}_S\equiv\bra{t}_C\ket{\Psi}$.
Since $\ket{\Psi}$ is the time-independent WdW state, the time derivative passes through as:
\begin{equation}\label{eq:chain}
  i\frac{d}{dt}\ket{\psi(t)}_S
    = i\frac{d}{dt}\bigl(\bra{t}_C\ket{\Psi}\bigr)
    = \Bigl(i\frac{d}{dt}\bra{t}_C\Bigr)\ket{\Psi}
    \stackrel{\eqref{eq:idbra}}{=} -\bra{t}_C\hat{H}_C\ket{\Psi}.
\end{equation}
Rearranging~\eqref{eq:chain} gives $\bra{t}_C\hat{H}_C\ket{\Psi}=-i\tfrac{d}{dt}\ket{\psi(t)}_S$.
Substituting into~\eqref{eq:sep}:
\begin{equation}\label{eq:prefinal}
  -i\frac{d}{dt}\ket{\psi(t)}_S+\hat{H}_S\ket{\psi(t)}_S=0.
\end{equation}

\subsection*{Step 4: Final Derivation}
Rearranging~\eqref{eq:prefinal} yields the standard Schr\"odinger form:
\begin{equation}\label{eq:schrodinger}
  i\frac{d}{dt}\ket{\psi(t)}_S = \hat{H}_S\ket{\psi(t)}_S.
\end{equation}
This is the central result of the Page--Wootters mechanism: unitary time evolution
emerges from the timeless WdW constraint $\hat{H}_{tot}\ket{\Psi}=0$ via
clock--system entanglement.

% -----------------------------------------------------------------------
\section{Mathematical Consistency Check: The Pauli Limit}
% -----------------------------------------------------------------------

In finite-dimensional qubit systems, the spectrum of $\hat{H}_C$ is discrete.
Two important consequences follow.

\begin{lemma}[Time Resolution]\label{lem:tres}
For a clock Hamiltonian $\hat{H}_C$ with operator norm $\|\hat{H}_C\|$, the minimum
resolvable time interval satisfies
\[
  \Delta t \;\ge\; \frac{1}{2\|\hat{H}_C\|},
  \qquad\text{i.e.\ }\Delta t\sim 1/\|\hat{H}_C\|.
\]
\end{lemma}
\begin{proof}
The Mandelstam--Tamm energy-time uncertainty relation states that for any observable $\hat{A}$:
\[
  \Delta E_C\cdot\tau_A \;\ge\; \tfrac{1}{2}\bigl|\langle[\hat{H}_C,\hat{A}]\rangle\bigr|,
\]
where $\tau_A=\Delta A\,/\,|\langle\dot{A}\rangle|$ is the characteristic evolution time of
$\hat{A}$.
Taking $\hat{A}$ to be the projector onto the initial clock state $\ket{0}_C$,
$\tau_A$ equals the minimum time for the clock to become distinguishable from its initial value.
Since $\Delta E_C\le\|\hat{H}_C\|$, we obtain
$\tau_A\ge(2\|\hat{H}_C\|)^{-1}$, i.e.\ $\Delta t\equiv\tau_A\sim 1/\|\hat{H}_C\|$.
\end{proof}

\begin{lemma}[Poincar\'e Recurrence]\label{lem:recurrence}
The relative state $\ket{\psi(t)}_S$ satisfying~\eqref{eq:schrodinger} is quasi-periodic
with recurrence time $T\le 2\pi/\Delta E_{min}$, where
$\Delta E_{min}:=\min_{k\neq j}|E_k^S-E_j^S|$.
\end{lemma}
\begin{proof}
Expand in the energy eigenbasis of $\hat{H}_S$:
\[
  \ket{\psi(t)}_S = \sum_j a_j\,e^{-iE_j^S t}\ket{f_j}_S, \qquad
  a_j=\braket{f_j|\psi(0)}_S.
\]
This is a finite sum of oscillatory terms (spectrum discrete by assumption).
The state returns near its initial value when all relative phases satisfy
$e^{-i(E_k^S-E_j^S)T}\approx 1$ simultaneously.
The slowest cycle is governed by $\Delta E_{min}$, requiring
$\Delta E_{min}\cdot T=2\pi$, hence $T\le 2\pi/\Delta E_{min}$.
\end{proof}

\begin{enumerate}
  \item \textbf{Discretization:} By Lemma~\ref{lem:tres}, time $t$ is sampled at resolutions
        $\Delta t\sim 1/\|\hat{H}_C\|$.
  \item \textbf{Poincar\'e Recurrence:} By Lemma~\ref{lem:recurrence},
        the state $\ket{\psi(t)}_S$ returns to its initial value after $T=2\pi/\Delta E_{min}$.
\end{enumerate}
This `error' is not a logical flaw but a physical consequence of finite information density
(Bekenstein bound) in qubit-based universes.

% -----------------------------------------------------------------------
\section{Noise as Graph Perturbation}
% -----------------------------------------------------------------------

In the BI-QMN framework, the system is represented by a weighted graph $G=(V,E,W)$, where
$V$ is the set of qubits and $W_{ij}$ represents the mutual information (or entanglement
measure) between qubits $i$ and $j$.

\begin{definition}[Graph Weight Function]\label{def:weight}
Given the $N$-qubit density matrix $\rho$, define the edge weight between qubits $i\neq j$
as $W_{ij}=|\rho_{ij}|$ (magnitude of off-diagonal coherence).
The degree matrix is $D_{ii}=\sum_{j\neq i}W_{ij}$, and the Graph Laplacian is
\begin{equation}\label{eq:lap}
  L = D-W
    = \sum_{(i,j)\in E}W_{ij}(e_i-e_j)(e_i-e_j)^T,
\end{equation}
where $e_k\in\mathbb{R}^{|V|}$ is the $k$-th standard basis vector.
\end{definition}

Under a Lindblad decoherence channel, off-diagonal elements decay as
$\rho_{ij}(t)=\rho_{ij}(0)e^{-\gamma_{ij}t}$, so for a time step $\Delta t$:
\[
  \delta W_{ij}
    = W_{ij}(0)\bigl(1-e^{-\gamma_{ij}\Delta t}\bigr)
    \approx W_{ij}(0)\,\gamma_{ij}\,\Delta t > 0.
\]
Each weight reduction contributes the rank-one perturbation
$-\delta W_{ij}(e_i-e_j)(e_i-e_j)^T$ to $L$, giving
\begin{equation}\label{eq:Lnoise}
  L_{noise} = L_0 + \Delta L, \qquad
  \Delta L = -\sum_{(i,j)\in E}\delta W_{ij}(e_i-e_j)(e_i-e_j)^T.
\end{equation}
Since each rank-one summand is positive semidefinite and $\delta W_{ij}>0$,
$\Delta L$ is \emph{negative semidefinite} ($\Delta L\preceq 0$).

% -----------------------------------------------------------------------
\section{Weyl's Inequality and Spectral Shift}
% -----------------------------------------------------------------------

\begin{theorem}[Weyl's Inequality]\label{thm:weyl}
Let $A$ and $B$ be $n\times n$ Hermitian matrices with ordered eigenvalues.  Then for all
$i=1,\dots,n$:
\[
  \lambda_i(A)+\lambda_{\min}(B)
    \;\le\; \lambda_i(A+B)
    \;\le\; \lambda_i(A)+\lambda_{\max}(B).
\]
\end{theorem}

Applying Theorem~\ref{thm:weyl} with $A=L_0$ and $B=\Delta L$, the perturbed eigenvalues
$\tilde\lambda_i$ satisfy:
\begin{equation}\label{eq:weyl}
  \lambda_i(L_0)+\lambda_{\min}(\Delta L)
    \;\le\;\tilde\lambda_i
    \;\le\;\lambda_i(L_0)+\lambda_{\max}(\Delta L).
\end{equation}
Because $\Delta L\preceq 0$, we have $\lambda_{\max}(\Delta L)\le 0$, so the upper bound gives
$\tilde\lambda_i\le\lambda_i(L_0)$: \emph{all eigenvalues shift downward}.
Furthermore $\mathrm{tr}(\Delta L)=-\sum_{i,j}\delta W_{ij}<0$ for any non-zero noise,
which forces $\lambda_{\min}(\Delta L)<0$.
The magnitude of the shift is bounded below by
\[
  |\lambda_{\min}(\Delta L)| = \|\Delta L\|_2
    \ge \frac{1}{\sqrt{n}}\|\Delta L\|_F > 0,
\]
where we used $\|B\|_2\ge\|B\|_F/\sqrt{n}$ for Hermitian $B$.
This confirms that any non-zero noise $\Delta L$ induces a deterministic, detectable shift in
the spectral density, providing a mathematical basis for anomaly detection.

% -----------------------------------------------------------------------
\section{Error Differentiation: $X$ vs.\ $Z$ Errors}
% -----------------------------------------------------------------------

We hypothesize that different quantum error channels map to distinct types of structural
distortions.

\subsection{Bit-Flip ($X$) Errors as Edge Deletion}

An $X$-error on qubit $k$ reduces the effective entanglement with its neighbors, decreasing
weights $\{W_{kj}\}_{j\sim k}$ by $\{\delta W_{kj}\}$.

\begin{proposition}[First-Order Spectral Shift under $X$-Errors]\label{prop:xshift}
Under an $X$-error that reduces edge weights $\{W_{kj}\}$ by $\{\delta W_{kj}>0\}$, the
first-order shift of the $m$-th Laplacian eigenvalue is
\begin{equation}\label{eq:xshift}
  \delta\lambda_m
    = -\sum_{j\sim k}\delta W_{kj}\bigl((v_m)_k-(v_m)_j\bigr)^2
    \;\le\; 0,
\end{equation}
where $\mathbf{v}_m$ is the $m$-th normalised eigenvector of $L_0$ (assuming non-degeneracy).
\end{proposition}

\begin{proof}
From Eq.~\eqref{eq:lap}, edge $(k,j)$ contributes $W_{kj}(e_k-e_j)(e_k-e_j)^T$ to $L$.
A weight reduction by $\delta W_{kj}$ induces the elementary perturbation
$\Delta L_{(k,j)}=-\delta W_{kj}(e_k-e_j)(e_k-e_j)^T$.
By first-order Rayleigh--Schr\"odinger perturbation theory for non-degenerate eigenvalues:
\begin{align*}
  \delta\lambda_m
    &= \mathbf{v}_m^T\Delta L\,\mathbf{v}_m
     = -\sum_{j\sim k}\delta W_{kj}\,
         \mathbf{v}_m^T(e_k-e_j)(e_k-e_j)^T\mathbf{v}_m \\
    &= -\sum_{j\sim k}\delta W_{kj}\bigl[(v_m)_k-(v_m)_j\bigr]^2.
\end{align*}
Non-negativity of the squared terms and positivity of $\delta W_{kj}$ ensure
$\delta\lambda_m\le 0$.
\end{proof}

For high-frequency eigenvectors $\mathbf{v}_m$ (large $m$), rapid spatial oscillation ensures
$|(v_m)_k-(v_m)_j|=\mathcal{O}(1)$, so the aggregate shift across all error edges is:
\begin{equation}\label{eq:xlarge}
  \Delta\lambda_{large} \;\propto\; -\sum_{(i,j)\in\mathrm{error}}\delta W_{ij}.
\end{equation}

\subsection{Phase-Flip ($Z$) Errors as Topological Warp}

$Z$-errors disrupt the phase relationship $\bra{t}_C\ket{\Psi}$, causing global connectivity
loss.  This is captured by the \emph{Fiedler value} $\lambda_2$, which bounds the edge
expansion of $G$ via the Cheeger inequality:
\[
  \frac{\lambda_2}{2} \;\le\; h(G) \;\le\; \sqrt{2\lambda_2\,\lambda_{\max}(L)},
\]
where $h(G)$ is the Cheeger constant.
The perturbed Fiedler value is given by the Courant--Fischer min-max theorem:
\begin{equation}\label{eq:fiedler}
  \tilde\lambda_2
    = \min_{\substack{x\perp\mathbf{1}\\\|x\|=1}}
      \frac{x^T(L_0+\Delta L)\,x}{x^Tx},
\end{equation}
where $\mathbf{1}\in\mathbb{R}^n$ denotes the all-ones vector.
Since $\Delta L\preceq 0$, we have $x^T\Delta L\,x\le 0$ for all $x$, therefore
\[
  \tilde\lambda_2
    \le \min_{\substack{x\perp\mathbf{1}\\\|x\|=1}} x^TL_0\,x = \lambda_2.
\]
The gap $\lambda_2-\tilde\lambda_2\ge 0$ directly quantifies the loss of algebraic connectivity
(global phase coherence).
A significant drop indicates \emph{Phase Collapse}.

\section{Relational Spectral Detection Without External Time}

\subsection{Clock-Conditioned Relative States}

We do not assume an external background time parameter. Instead, following the Page--Wootters
construction, we consider a timeless global state
$|\Psi\rangle \in \mathcal{H}_C \otimes \mathcal{H}_S$ satisfying the global constraint
\begin{equation}
  \hat{H}_{\mathrm{tot}} |\Psi\rangle = 0.
\end{equation}

Let $\{|\tau\rangle_C\}$ denote the clock basis. The corresponding clock-conditioned relative
state is defined as
\begin{equation}
  |\psi(\tau)\rangle_S := {}_C\langle \tau | \Psi \rangle,
\end{equation}
or, in density-matrix form,
\begin{equation}
  \rho_S(\tau)
    :=
    \frac{
      \mathrm{Tr}_C \!\left[
        \left(|\tau\rangle\langle\tau|_C \otimes I_S\right)
        |\Psi\rangle\langle\Psi|
      \right]
    }{
      \mathrm{Tr}\!\left[
        \left(|\tau\rangle\langle\tau|_C \otimes I_S\right)
        |\Psi\rangle\langle\Psi|
      \right]
    }.
\end{equation}

Here, $\tau$ is not an external time variable, but an internal relational label induced by the
clock subsystem.

\subsection{Relational Graph and Laplacian Family}

For each clock-conditioned state $\rho_S(\tau)$, we construct a weighted graph
\begin{equation}
  G_{\tau} = (V, E, W^{(\tau)}),
\end{equation}
where the edge weights are determined from the relational structure of $\rho_S(\tau)$.

For example, one may define
\begin{equation}
  W^{(\tau)}_{ij} := f\!\left(\rho_S(\tau); i,j\right),
\end{equation}
where $f$ denotes a pairwise relational measure such as coherence magnitude, mutual
information, or another correlation functional.

The corresponding degree matrix is
\begin{equation}
  D^{(\tau)}_{ii} = \sum_{j \neq i} W^{(\tau)}_{ij},
\end{equation}
and the clock-conditioned graph Laplacian is
\begin{equation}
  L(\tau) = D^{(\tau)} - W^{(\tau)}.
\end{equation}

Thus, instead of a time-evolving Laplacian, we obtain a relationally indexed family of
Laplacians
\begin{equation}
  \{L(\tau)\}_{\tau \in \mathcal{T}},
\end{equation}
where $\mathcal{T}$ denotes the clock-label domain.

\subsection{Relational Spectral Density}

Let
\begin{equation}
  \Lambda(\tau) = \{\lambda_1(\tau), \dots, \lambda_n(\tau)\}
\end{equation}
be the eigenvalue set of $L(\tau)$.

We define the clock-conditioned empirical spectral density (ESD) as
\begin{equation}
  p_{\tau}(\lambda)
    = \frac{1}{n} \sum_{k=1}^{n} \delta\!\bigl(\lambda - \lambda_k(\tau)\bigr),
\end{equation}
or, numerically, via a kernel density estimator,
\begin{equation}
  p_{\tau}(\lambda)
    = \frac{1}{n} \sum_{k=1}^{n}
      \frac{1}{\sqrt{2\pi \sigma_h^2}}
      \exp\!\left(-\frac{(\lambda-\lambda_k(\tau))^2}{2\sigma_h^2}\right).
\end{equation}

Accordingly, spectral variation is interpreted as variation across clock-conditioned relational
slices, rather than evolution over an external time axis.

\subsection{Reference and Observed Relational Spectra}

For each relational slice $\tau$, let $p_{\tau}^{\mathrm{ref}}(\lambda)$ denote the reference
spectral density and $p_{\tau}^{\mathrm{obs}}(\lambda)$ the observed spectral density under
perturbation.

The slice-wise information-theoretic discrepancy is defined as
\begin{equation}
  D_{\mathrm{KL}}^{(\tau)}
    := D_{\mathrm{KL}}
       \!\left(p_{\tau}^{\mathrm{ref}} \,\middle\|\, p_{\tau}^{\mathrm{obs}}\right).
\end{equation}

\begin{theorem}[Relational Spectral Detectability]
For a fixed clock-conditioned slice $\tau$, if the perturbation induces a non-trivial spectral
deformation, $\|\Delta L(\tau)\|_F > 0$, then $D_{\mathrm{KL}}^{(\tau)} > 0$.
\end{theorem}

\begin{proof}
If $\|\Delta L(\tau)\|_F > 0$, then by Weyl's inequality (Theorem~\ref{thm:weyl}) at least one
eigenvalue of $L(\tau)$ is shifted relative to the reference Laplacian.  Hence
$p_{\tau}^{\mathrm{ref}}(\lambda) \neq p_{\tau}^{\mathrm{obs}}(\lambda)$ on a set of non-zero
measure.  By Gibbs' inequality,
$D_{\mathrm{KL}}\!\left(p_{\tau}^{\mathrm{ref}} \,\middle\|\, p_{\tau}^{\mathrm{obs}}\right)>0$.
\end{proof}

Thus, $D_{\mathrm{KL}}^{(\tau)}$ serves as a sufficient certificate of non-trivial spectral
deformation within the relational slice indexed by $\tau$.

\subsection{Relationally Averaged Detection Functional}

To compare the entire relational family rather than a single slice, we define the aggregated
detection functional
\begin{equation}
  \mathcal{D}_{\mathrm{KL}}
    = \int_{\mathcal{T}}
      D_{\mathrm{KL}}\!\left(p_{\tau}^{\mathrm{ref}} \,\middle\|\, p_{\tau}^{\mathrm{obs}}\right)
      d\mu(\tau),
\end{equation}
where $\mu(\tau)$ is a measure over the clock-label domain $\mathcal{T}$.

In discrete simulations, this reduces to
\begin{equation}
  \mathcal{D}_{\mathrm{KL}}
    = \frac{1}{|\mathcal{T}|} \sum_{\tau \in \mathcal{T}} D_{\mathrm{KL}}^{(\tau)}.
\end{equation}

This quantity measures the global mismatch between two relational spectral families without
introducing an external time parameter.

\subsection{Hybrid Detection Metric with Geometric Augmentation}

While $D_{\mathrm{KL}}^{(\tau)}$ provides an information-theoretic detection certificate, it
does not directly quantify the geometric magnitude of spectral displacement.  To incorporate
geometric sensitivity, we define
\begin{equation}
  D_{\mathrm{detect}}^{(\tau)}
    = \alpha\,\widetilde{D}_{\mathrm{KL}}^{(\tau)}
      + (1-\alpha)\,\widetilde{W}_1^{(\tau)},
\end{equation}
where
\begin{equation}
  \widetilde{D}_{\mathrm{KL}}^{(\tau)}
    = \mathrm{Norm}\!\left[
        D_{\mathrm{KL}}\!\left(p_{\tau}^{\mathrm{ref}} \,\middle\|\, p_{\tau}^{\mathrm{obs}}\right)
      \right]
\end{equation}
and
\begin{equation}
  \widetilde{W}_1^{(\tau)}
    = \mathrm{Norm}\!\left[
        W_1\!\left(p_{\tau}^{\mathrm{ref}}, p_{\tau}^{\mathrm{obs}}\right)
      \right].
\end{equation}

Here, $W_1$ denotes the Wasserstein-1 distance between the two clock-conditioned spectral
densities, and $\alpha \in [0,1]$ controls the balance between information-theoretic
sensitivity and geometric displacement.

We then define the relationally aggregated hybrid score
\begin{equation}
  \mathcal{D}_{\mathrm{detect}}
    = \int_{\mathcal{T}} D_{\mathrm{detect}}^{(\tau)}\,d\mu(\tau),
\end{equation}
or, in discrete form,
\begin{equation}
  \mathcal{D}_{\mathrm{detect}}
    = \frac{1}{|\mathcal{T}|} \sum_{\tau \in \mathcal{T}} D_{\mathrm{detect}}^{(\tau)}.
\end{equation}

\begin{remark}
The KL term preserves the information-theoretic detectability guarantee at each relational
slice, while the Wasserstein term quantifies geometric spectral displacement across the
clock-conditioned family.
\end{remark}

\subsection{Symmetric Stability Check}

For numerical robustness, we optionally introduce the Jensen--Shannon divergence at each
relational slice:
\begin{equation}
  D_{\mathrm{JS}}^{(\tau)}
    = \frac{1}{2}
      D_{\mathrm{KL}}\!\left(p_{\tau}^{\mathrm{ref}} \,\middle\|\, m_{\tau}\right)
      + \frac{1}{2}
      D_{\mathrm{KL}}\!\left(p_{\tau}^{\mathrm{obs}} \,\middle\|\, m_{\tau}\right),
\end{equation}
with
\begin{equation}
  m_{\tau} = \frac{1}{2}\left(p_{\tau}^{\mathrm{ref}} + p_{\tau}^{\mathrm{obs}}\right).
\end{equation}

The relationally averaged Jensen--Shannon score is
\begin{equation}
  \mathcal{D}_{\mathrm{JS}}
    = \int_{\mathcal{T}} D_{\mathrm{JS}}^{(\tau)}\,d\mu(\tau),
\end{equation}
or its discrete counterpart.  This quantity is used as a symmetric numerical stability check,
rather than as the primary detectability certificate.

\subsection{Localized Geometric Error Differentiation}

To distinguish different structural perturbation types, we introduce a band-limited Wasserstein
descriptor on a spectral region $\mathcal{B}$:
\begin{equation}
  W_1^{(\tau,\mathcal{B})}
    = \int_{\mathcal{B}}
      \left|F_{\tau}^{\mathrm{ref}}(\lambda) - F_{\tau}^{\mathrm{obs}}(\lambda)\right|
      d\lambda,
\end{equation}
where $F_{\tau}^{\mathrm{ref}}$ and $F_{\tau}^{\mathrm{obs}}$ are the corresponding cumulative
spectral distributions.

This enables local-versus-global deformation analysis within each relational slice:
\begin{itemize}
  \item high-frequency bands capture localized structural perturbations,
  \item low-frequency bands capture global connectivity loss.
\end{itemize}

Thus, error discrimination is formulated as a geometric comparison across relational spectral
slices rather than as a trajectory over background time.

\subsection{Detection Criterion}

An anomaly is declared whenever $\mathcal{D}_{\mathrm{detect}} > \varepsilon$, where
$\varepsilon$ is determined from the background distribution of relationally aggregated
divergence values under unperturbed fluctuations:
\begin{equation}
  \varepsilon
    = \mathbb{E}[\mathcal{D}_{\mathrm{vac}}]
      + \kappa\,\mathrm{Std}[\mathcal{D}_{\mathrm{vac}}].
\end{equation}

This criterion remains fully relational and does not require any external time indexing.

\section{Relational Spectral Trajectory Manifold}

\subsection{Clock-Conditioned Spectral Family}

Let $|\Psi\rangle \in \mathcal{H}_C \otimes \mathcal{H}_S$ be a timeless global state satisfying
\begin{equation}
\hat{H}_{\mathrm{tot}} |\Psi\rangle = 0.
\end{equation}

For each clock label $\tau \in \mathcal{T}$, define the clock-conditioned relative state
\begin{equation}
\rho_S(\tau)
:=
\frac{
\mathrm{Tr}_C \!\left[
\left(|\tau\rangle\langle\tau|_C \otimes I_S\right)
|\Psi\rangle\langle\Psi|
\right]
}{
\mathrm{Tr}\!\left[
\left(|\tau\rangle\langle\tau|_C \otimes I_S\right)
|\Psi\rangle\langle\Psi|
\right]
}.
\end{equation}

From each relational slice $\rho_S(\tau)$, construct the weighted graph Laplacian
\begin{equation}
L(\tau) = D^{(\tau)} - W^{(\tau)},
\end{equation}
and let its ordered eigenvalue vector be
\begin{equation}
\Lambda(\tau) = \bigl(\lambda_1(\tau), \lambda_2(\tau), \dots, \lambda_n(\tau)\bigr)
\in \mathbb{R}^n.
\end{equation}

We define the \emph{relational spectral family} associated with $|\Psi\rangle$ as
\begin{equation}
\Gamma_{\Psi} : \mathcal{T} \to \mathbb{R}^n,
\qquad
\Gamma_{\Psi}(\tau) := \Lambda(\tau).
\end{equation}

Thus, $\Gamma_{\Psi}$ is not a trajectory over external time, but an ordered family of spectral signatures indexed by the internal clock labels.

\subsection{Spectral Trajectory Space}

Let $\mathcal{M}_{\mathrm{spec}}$ denote the set of admissible relational spectral families,
\begin{equation}
\mathcal{M}_{\mathrm{spec}}
:=
\left\{
\Gamma : \mathcal{T} \to \mathbb{R}^n
\ \middle|\
\Gamma(\tau) = \Lambda\!\left(L(\rho_S(\tau))\right)
\text{ for some physically admissible } \rho_S(\tau)
\right\}.
\end{equation}

We interpret $\mathcal{M}_{\mathrm{spec}}$ as the \emph{relational spectral trajectory manifold}, i.e., the manifold of all physically realizable clock-conditioned spectral families.

\begin{remark}
Each element of $\mathcal{M}_{\mathrm{spec}}$ is a full relational trajectory $\Gamma$, not a single spectral snapshot.
\end{remark}
\begin{remark}
Strictly speaking, $\mathcal{M}_{\mathrm{spec}}$ is treated here 
as a metric trajectory space of admissible spectral families 
rather than as a fully established smooth manifold. The geodesic 
formulation in Secs.~9.3--9.9 should therefore be interpreted as 
a formal geometric model. The discrete representation in Eq.~(84) 
provides the computational realization, in which 
$\mathcal{M}_{\mathrm{spec}}$ reduces to a constrained subset of 
$\mathbb{R}^{mn}$.
\end{remark}

\subsection{Tangent Structure}

Given a smooth one-parameter family of relational trajectories
\begin{equation}
\Gamma_s \in \mathcal{M}_{\mathrm{spec}}, \qquad s \in (-\epsilon,\epsilon),
\end{equation}
with $\Gamma_0 = \Gamma$, the tangent vector at $\Gamma$ is defined by
\begin{equation}
V(\tau)
=
\left.\frac{d}{ds}\Gamma_s(\tau)\right|_{s=0}
\in \mathbb{R}^n.
\end{equation}

Hence, the tangent space at $\Gamma$ is
\begin{equation}
T_{\Gamma}\mathcal{M}_{\mathrm{spec}}
=
\left\{
V : \mathcal{T} \to \mathbb{R}^n
\ \middle|\
V(\tau)=\left.\frac{d}{ds}\Gamma_s(\tau)\right|_{s=0}
\text{ for some smooth family } \Gamma_s
\right\}.
\end{equation}

Intuitively, $V$ describes an infinitesimal deformation of the full relational spectral family.

\subsection{Riemannian Metric on the Trajectory Manifold}

To quantify distances between relational trajectories, we introduce a Riemannian metric on $\mathcal{M}_{\mathrm{spec}}$.

For tangent vectors $U, V \in T_{\Gamma}\mathcal{M}_{\mathrm{spec}}$, define
\begin{equation}
g_{\Gamma}(U,V)
=
\int_{\mathcal{T}}
U(\tau)^{\top}
\, G\!\bigl(\Gamma(\tau)\bigr)\,
V(\tau)
\, d\mu(\tau),
\end{equation}
where $G(\Gamma(\tau))$ is a positive-definite local metric tensor on $\mathbb{R}^n$, and $\mu$ is a measure over the clock-label domain $\mathcal{T}$.

A simple Euclidean choice is
\begin{equation}
G(\Gamma(\tau)) = I_n,
\end{equation}
which yields
\begin{equation}
g_{\Gamma}(U,V)
=
\int_{\mathcal{T}}
U(\tau)^{\top} V(\tau)\, d\mu(\tau).
\end{equation}

More generally, one may choose $G$ to weight low-frequency and high-frequency spectral components differently.

\subsection{Trajectory Length and Distance}

Let $\gamma : [0,1] \to \mathcal{M}_{\mathrm{spec}}$ be a smooth path in the trajectory manifold, where each $\gamma(s)$ is itself a relational spectral family:
\begin{equation}
\gamma(s) = \Gamma_s.
\end{equation}

Its velocity field is
\begin{equation}
\dot{\gamma}(s) = \frac{d}{ds}\Gamma_s \in T_{\Gamma_s}\mathcal{M}_{\mathrm{spec}}.
\end{equation}

The length of $\gamma$ is defined by
\begin{equation}
\mathrm{Len}(\gamma)
=
\int_0^1
\sqrt{
g_{\Gamma_s}\!\bigl(\dot{\gamma}(s),\dot{\gamma}(s)\bigr)
}
\, ds.
\end{equation}

The induced geodesic distance between two relational trajectories
$\Gamma^{(a)}, \Gamma^{(b)} \in \mathcal{M}_{\mathrm{spec}}$
is
\begin{equation}
d_{\mathcal{M}}
\!\left(
\Gamma^{(a)}, \Gamma^{(b)}
\right)
=
\inf_{\gamma(0)=\Gamma^{(a)},\, \gamma(1)=\Gamma^{(b)}}
\mathrm{Len}(\gamma).
\end{equation}

\subsection{Observed and Reference Trajectories}

Let
\begin{equation}
\Gamma_{\mathrm{ref}} \in \mathcal{M}_{\mathrm{spec}}
\end{equation}
denote the reference relational spectral trajectory, and let
\begin{equation}
\Gamma_{\mathrm{obs}} \in \mathcal{M}_{\mathrm{spec}}
\end{equation}
denote the observed trajectory under perturbation.

The trajectory-level distortion is then defined as
\begin{equation}
\Delta_{\mathrm{traj}}
=
d_{\mathcal{M}}
\!\left(
\Gamma_{\mathrm{ref}}, \Gamma_{\mathrm{obs}}
\right).
\end{equation}

This quantity measures the global deformation of the entire clock-conditioned spectral family, rather than the discrepancy at any single relational slice.

\subsection{Geodesic Recovery Principle}

We now formulate recovery as a geometric restoration problem on $\mathcal{M}_{\mathrm{spec}}$.

Given an observed trajectory $\Gamma_{\mathrm{obs}}$, define the recovered trajectory as
\begin{equation}
\Gamma^{*}
=
\arg\min_{\Gamma \in \mathcal{A}}
d_{\mathcal{M}}
\!\left(
\Gamma, \Gamma_{\mathrm{obs}}
\right),
\end{equation}
where $\mathcal{A} \subset \mathcal{M}_{\mathrm{spec}}$ is the admissible set of physically consistent or reference-compatible trajectories.

Equivalently, if $\Gamma_{\mathrm{ref}}$ is known, the recovery path is the minimizing geodesic
\begin{equation}
\gamma^{*}
=
\arg\min_{\gamma(0)=\Gamma_{\mathrm{obs}},\,\gamma(1)=\Gamma_{\mathrm{ref}}}
\mathrm{Len}(\gamma).
\end{equation}

This gives the central interpretation:
\begin{quote}
Recovery is not modeled as evolution over background time, but as the shortest admissible path on the relational spectral trajectory manifold.
\end{quote}

\subsection{Geodesic Equation}

Let the path be locally represented by coordinates
\begin{equation}
x^a(s), \qquad a=1,\dots,\dim(\mathcal{M}_{\mathrm{spec}}),
\end{equation}
with Riemannian metric $g_{ab}(x)$. The geodesic equation is
\begin{equation}
\frac{d^2 x^a}{ds^2}
+
\Gamma^{a}_{bc}(x)
\frac{dx^b}{ds}
\frac{dx^c}{ds}
= 0,
\end{equation}
where the Christoffel symbols are
\begin{equation}
\Gamma^{a}_{bc}
=
\frac{1}{2}
g^{ad}
\left(
\partial_b g_{dc}
+
\partial_c g_{db}
-
\partial_d g_{bc}
\right).
\end{equation}

In the Euclidean metric case $G=I_n$, the connection is flat and the geodesic reduces to a straight interpolation in trajectory space:
\begin{equation}
\Gamma_s(\tau)
=
(1-s)\Gamma_{\mathrm{obs}}(\tau)
+
s\,\Gamma_{\mathrm{ref}}(\tau),
\qquad s \in [0,1].
\end{equation}

\subsection{Variational Formulation}

Equivalently, geodesics are stationary points of the action
\begin{equation}
\mathcal{S}[\gamma]
=
\frac{1}{2}
\int_0^1
g_{\Gamma_s}\!\bigl(\dot{\gamma}(s),\dot{\gamma}(s)\bigr)\, ds.
\end{equation}

Thus, the geodesic recovery problem can be formulated as
\begin{equation}
\gamma^{*}
=
\arg\min_{\gamma}
\mathcal{S}[\gamma],
\end{equation}
subject to the endpoint conditions
\begin{equation}
\gamma(0)=\Gamma_{\mathrm{obs}},
\qquad
\gamma(1)=\Gamma_{\mathrm{ref}},
\end{equation}
or more generally,
\begin{equation}
\gamma(1)\in\mathcal{A}.
\end{equation}

\section{Admissible Relational Spectral Trajectory Set}

\subsection{Definition of the Admissible Set}

Let $\mathcal{M}_{\mathrm{spec}}$ denote the relational spectral trajectory manifold.
We define the admissible set $\mathcal{A} \subset \mathcal{M}_{\mathrm{spec}}$ as
\begin{equation}
\mathcal{A}
=
\mathcal{A}_{\mathrm{lap}}
\cap
\mathcal{A}_{\mathrm{rel}}
\cap
\mathcal{A}_{\mathrm{ref}},
\end{equation}
where each component enforces a distinct class of physical and structural constraints.

\subsection{Laplacian Consistency}

A trajectory $\Gamma \in \mathcal{M}_{\mathrm{spec}}$ belongs to $\mathcal{A}_{\mathrm{lap}}$
if for all $\tau \in \mathcal{T}$, the associated matrix $L(\tau)$ satisfies

\begin{align}
L(\tau) &= L(\tau)^\top, \\
L(\tau) &\succeq 0, \\
L(\tau)\mathbf{1} &= 0, \\
L_{ij}(\tau) &\le 0 \quad (i \neq j), \\
L_{ii}(\tau) &= \sum_{j\neq i} -L_{ij}(\tau).
\end{align}

Thus,
\begin{equation}
\mathcal{A}_{\mathrm{lap}}
=
\left\{
\Gamma \in \mathcal{M}_{\mathrm{spec}}
\;\middle|\;
L(\tau)\ \text{is a valid graph Laplacian for all }\tau
\right\}.
\end{equation}

\subsection{Relational Consistency}

A trajectory $\Gamma$ is relationally admissible if it satisfies

\paragraph{(i) Slice Regularity}
\begin{equation}
\int_{\mathcal{T}}
\left\|
\frac{d}{d\tau}\Gamma(\tau)
\right\|^2
\, d\mu(\tau)
< \infty.
\end{equation}

\paragraph{(ii) Spectral Ordering}
\begin{equation}
\lambda_1(\tau) \le \lambda_2(\tau) \le \cdots \le \lambda_n(\tau)
\quad \forall \tau \in \mathcal{T}.
\end{equation}

\paragraph{(iii) Clock Consistency}
Let $d_{\mathrm{clock}}(\rho_S(\tau))$ be a clock-consistency measure. Then
\begin{equation}
\int_{\mathcal{T}}
d_{\mathrm{clock}}(\rho_S(\tau))
\, d\mu(\tau)
\le \eta.
\end{equation}

Thus,
\begin{equation}
\mathcal{A}_{\mathrm{rel}}
=
\left\{
\Gamma \in \mathcal{A}_{\mathrm{lap}}
\;\middle|\;
\Gamma\ \text{is slice-regular and clock-consistent}
\right\}.
\end{equation}

\subsection{Reference Compatibility}

Let $\mathfrak{R} \subset \mathcal{M}_{\mathrm{spec}}$ be a reference trajectory family.
We define

\begin{equation}
\mathcal{A}_{\mathrm{ref}}
=
\left\{
\Gamma \in \mathcal{A}_{\mathrm{rel}}
\;\middle|\;
\inf_{\Gamma' \in \mathfrak{R}}
d_{\mathcal{M}}(\Gamma,\Gamma')
\le \delta
\right\}.
\end{equation}

\subsection{Final Definition}

The admissible set is given by
\begin{equation}
\mathcal{A}
=
\left\{
\Gamma \in \mathcal{M}_{\mathrm{spec}}
\;\middle|\;
\begin{array}{l}
L(\tau)\ \text{is a valid Laplacian for all }\tau,\\[4pt]
\Gamma\ \text{is slice-regular and clock-consistent},\\[4pt]
\inf_{\Gamma' \in \mathfrak R} d_{\mathcal M}(\Gamma,\Gamma') \le \delta
\end{array}
\right\}.
\end{equation}

\subsection{Discrete Computational Form}

Let $\mathcal{T} = \{\tau_1,\dots,\tau_m\}$. Then a trajectory is represented as
\begin{equation}
\mathbf{z}_\Gamma
=
\begin{bmatrix}
\Lambda(\tau_1) \\
\Lambda(\tau_2) \\
\vdots \\
\Lambda(\tau_m)
\end{bmatrix}
\in \mathbb{R}^{mn}.
\end{equation}

The trajectory distance becomes
\begin{equation}
d_{\mathcal{M}}(\Gamma^{(a)},\Gamma^{(b)})
=
\left(
\sum_{r=1}^m
w_r
\left\|
\Lambda^{(a)}(\tau_r)-\Lambda^{(b)}(\tau_r)
\right\|_2^2
\right)^{1/2}.
\end{equation}

\subsection{Penalty-Based Admissibility}

Define

\begin{align}
\Phi_{\mathrm{lap}}(\Gamma)
&=
\sum_{r=1}^m
\left(
\|L(\tau_r)-L(\tau_r)^\top\|_F^2
+
\|\min(0,\lambda(L(\tau_r)))\|_2^2
+
\|L(\tau_r)\mathbf{1}\|_2^2
\right), \\[6pt]
\Phi_{\mathrm{smooth}}(\Gamma)
&=
\sum_{r=1}^{m-1}
\frac{
\|\Lambda(\tau_{r+1})-\Lambda(\tau_r)\|_2^2
}{
(\tau_{r+1}-\tau_r)^2
}, \\[6pt]
\Phi_{\mathrm{ref}}(\Gamma)
&=
\min_{\Gamma' \in \mathfrak R}
d_{\mathcal M}(\Gamma,\Gamma')^2, \\[6pt]
\Phi_{\mathrm{clock}}^{\mathrm{geom}}(\Gamma)&=\sum_{r=1}^{m-1} \frac{\left\|\Lambda\left(\tau_{r+1}\right)-\Lambda\left(\tau_r\right)\right\|_2}{d_{\mathrm{clock}}^{\mathrm{geom}}\left(\tau_r, \tau_{r+1}\right)+\epsilon}
\end{align}

Here $\Phi_{\mathrm{lap}}$ is a graph-validity surrogate, $\Phi_{\mathrm{smooth}}$ is a regularizer on relational variation, and $\Phi_{\mathrm{clock}}$ is a relational consistency penalty.  The quantity $\Phi_{\mathrm{ref}}$ remains useful as a coarse reference-compatibility diagnostic, but it does not by itself provide a directed restoration force toward the clean family.

We therefore define the relaxed admissible set
\begin{equation}
\mathcal{A}_\epsilon
=
\left\{
\Gamma :
\Phi_{\mathrm{lap}} \le \epsilon_1,\;
\Phi_{\mathrm{smooth}} \le \epsilon_2,\;
\Phi_{\mathrm{ref}} \le \epsilon_3,\;
\Phi_{\mathrm{clock}} \le \epsilon_4
\right\}.
\end{equation}

\subsection{Recovery Objective}

The earlier recovery objective primarily enforced admissibility and observational consistency, but did not sufficiently encode restoration toward the reference family. To address this mismatch, we introduce an explicit reference-guided recovery term based on the nearest admissible reference anchor. This modification separates restoration from feasibility, allowing the optimization to favor trajectories that are both admissible and structurally closer to the clean reference family.

Let
\begin{equation}
\Gamma_{\mathrm{ref,anchor}}
:=
\arg\min_{\Gamma' \in \mathfrak R}
d_{\mathcal M}(\Gamma',\Gamma_{\mathrm{obs}})^2
\end{equation}
denote the nearest reference anchor to the observed trajectory.

The recovered trajectory is then defined by the restoration-aware objective
\begin{equation}
\Gamma^*
=
\arg\min_{\Gamma}
\Bigl[
\alpha\, d_{\mathcal M}(\Gamma,\Gamma_{\mathrm{obs}})^2
+
\beta\, d_{\mathcal M}(\Gamma,\Gamma_{\mathrm{ref,anchor}})^2
+
\lambda_1 \Phi_{\mathrm{lap}}(\Gamma)
+
\lambda_2 \Phi_{\mathrm{smooth}}(\Gamma)
+
\lambda_3 \Phi_{\mathrm{clock}}(\Gamma)
\Bigr].
\end{equation}

In this formulation, the observed-fit term preserves consistency with the measured trajectory, the anchor term supplies an explicit restoration direction toward the reference family, and the remaining penalties serve only as feasibility / regularization controls.

\subsection{Two-Stage Admissible Recovery}

We refine the recovery procedure by introducing a two-stage structure on the admissible trajectory space.

\paragraph{Stage 1 (Discrete Admissible Projection).}
Given an observed trajectory $\Gamma_{\mathrm{obs}}$, we first project onto a discrete admissible candidate bank $\mathcal{A}_{\mathrm{bank}} \subset \mathcal{A}$:
\[
\Gamma^{(1)} 
= 
\arg\min_{\Gamma \in \mathcal{A}_{\mathrm{bank}}}
\left[
\alpha\, d_M(\Gamma, \Gamma_{\mathrm{obs}})^2
+
\beta\, d_M(\Gamma, \Gamma_{\mathrm{ref,anchor}})^2
+
\lambda_1 \Phi_{\mathrm{lap}}(\Gamma)
+
\lambda_2 \Phi_{\mathrm{smooth}}(\Gamma)
+
\lambda_3 \Phi_{\mathrm{clock}}(\Gamma)
\right].
\]
This step provides a stable admissible anchor while explicitly preserving restoration toward the nearest compatible reference family.

\paragraph{Stage 2 (Admissible Convex Refinement).}
We optionally perform a continuous refinement within a local convex hull of admissible candidates.
Let $\{\Gamma_k\}$ denote a selected subset of admissible candidates. We consider convex combinations
\[
\Gamma(\alpha) = \sum_k \alpha_k \Gamma_k, \quad \alpha_k \ge 0, \quad \sum_k \alpha_k = 1,
\]
restricted to those $\alpha$ such that $\Gamma(\alpha) \in \mathcal{A}$.

The refined trajectory is defined as
\[
\Gamma^{(2)} 
=
\arg\min_{\Gamma(\alpha) \in \mathcal{A}}
\left[
\alpha\, d_M(\Gamma(\alpha), \Gamma_{\mathrm{obs}})^2
+
\beta\, d_M(\Gamma(\alpha), \Gamma_{\mathrm{ref,anchor}})^2
+
\lambda_1 \Phi_{\mathrm{lap}}(\Gamma(\alpha))
+
\lambda_2 \Phi_{\mathrm{smooth}}(\Gamma(\alpha))
+
\lambda_3 \Phi_{\mathrm{clock}}(\Gamma(\alpha))
\right].
\]

\paragraph{Interpretation.}
Stage 1 performs a coarse projection onto an admissible discrete set, while Stage 2 carries out a local refinement under the same restoration-aware objective.  The reference-anchor term is the actual recovery driver, whereas $\Phi_{\mathrm{lap}}$, $\Phi_{\mathrm{smooth}}$, and $\Phi_{\mathrm{clock}}$ retain their roles as feasibility / regularization terms.

\paragraph{Practical Use.}
Stage 1 is used as the default recovery due to its stability and interpretability. Stage 2 is treated as an optional diagnostic refinement and is adopted only when it decreases the restoration-aware objective while maintaining admissibility. In practice this means checking
\[
J_{\mathrm{rec}}(\Gamma^{(2)}) < J_{\mathrm{rec}}(\Gamma^{(1)})
\]
together with the required admissibility constraints, e.g.
\[
\Phi_{\mathrm{lap}}(\Gamma^{(2)}) \le \epsilon_1,
\qquad
\Phi_{\mathrm{clock}}(\Gamma^{(2)}) \le \epsilon_4.
\]

\paragraph{Default Operating Point.}
For the paper-facing default configuration, the discrete reference bank is taken to have width
\[
\Delta_{\mathrm{bank}} = 20^\circ
\]
around the central reference phase, and the restoration weight is set to
\[
\beta = 16.
\]
This operating point is chosen as the smallest anchor strength that already gives a stable Stage 1 surrogate in the validation sweep.  Narrower banks such as $5^\circ$ are treated as unstable narrow-bank regimes, while very large anchor weights such as $\beta = 128$ are treated as high-anchor regimes that overly suppress useful Stage 2 diagnostic variation.  Wider operating points or stronger anchors are therefore relegated to sensitivity / ablation tables rather than used as the main paper default.

\subsection{Hybrid Policy Layer for Encoded Baselines}

The encoded-QEC experiments produce two recovery candidates for each sampled
noise instance:
\begin{align*}
A &:= \text{syndrome-first decoded recovery},\\
B &:= \text{relational recovery from the admissible restoration pipeline}.
\end{align*}
Candidate $A$ represents the conventional decoder-facing baseline.  Candidate $B$
represents the relational restoration mechanism developed above: it is selected from
the admissible trajectory family using the restoration-aware objective rather than
from the syndrome prior alone.  The hybrid layer is therefore not a third recovery
algorithm.  It is a decision layer that asks when the relational candidate is reliable
enough to replace the syndrome candidate.

\subsubsection*{Candidate diagnostics}

For each candidate $X\in\{A,B\}$ we record:
\[
F_X,\qquad
\Delta F_X:=F_X-F_{\mathrm{before}},\qquad
J_X:=J_{\mathrm{rec}}(\Gamma_X),
\]
together with the admissibility report produced by the relational trajectory
constraints.  Let $\mathcal{K}$ be the set of active admissibility penalties, and let
$\theta_k$ denote the corresponding threshold for penalty $\Phi_k$.  The normalized
admissibility violation used by the robust gate is
\[
v_X
:=
\max_{k\in\mathcal{K}}
\max\!\left\{0,\frac{\Phi_k(\Gamma_X)-\theta_k}{\theta_k+\epsilon}\right\}.
\]
Thus $v_X=0$ means that all active constraints are within threshold, while positive
values measure the worst normalized excursion beyond the admissible boundary.  This
is strictly more informative than a Boolean admissible/inadmissible flag, and it
matches the implementation used by the robust gate.

\subsubsection*{Decision variables used by the hybrid layer}

The policies below are deterministic maps from the same row-level diagnostics to one
of the two candidates.  Let
\[
a_X:=\mathbf{1}[X\ \mathrm{is\ admissible}],
\qquad
d_X:=d(\Gamma_X,\Gamma_{\mathrm{obs}}),
\qquad
u_X:=\frac{1}{1+J_X}.
\]
The syndrome support score used in the current implementation is
\[
s_A:=p_{\mathrm{obs}}(0),
\qquad
s_B:=0,
\]
where $p_{\mathrm{obs}}(0)$ is the observed mean no-error syndrome probability.  This
encodes the prior that the decoder-facing candidate $A$ is directly supported by the
syndrome record, while the relational candidate $B$ must earn selection through
trajectory proximity, restoration objective, and admissibility.  Under the imperfect
syndrome protocols, the degraded syndrome variables enter this policy layer; the
physical construction of $A$ and $B$ is otherwise kept fixed so that the experiment
isolates decision risk rather than changing the recovery candidates themselves.

\subsubsection*{C1: rule-based conservative hybrid}

The first policy is the original conservative hybrid used as a reference point.  Let
$c_{\mathrm{obs}}$ denote observed syndrome consistency and let $\tau_J$ be the
objective tie tolerance.  In the seed-10 experiments
\[
\tau_J=10^{-9},
\qquad
c_{\mathrm{obs}}\ \mathrm{is\ required\ for\ the\ objective\ tie\ switch}.
\]
The implemented decision rule is
\[
C_1=B
\quad\Longleftrightarrow\quad
\bigl[(a_A=0)\wedge(a_B=1)\bigr]
\ \vee\
\bigl[(a_B=1)\wedge(J_B<J_A-\tau_J)\wedge c_{\mathrm{obs}}\bigr].
\]
Otherwise $C_1=A$.  Thus $C_1$ is not a fitted score model.  It has two explicit
routes to $B$: veto a non-admissible syndrome candidate when $B$ is admissible, or use
$B$ as an objective tie-breaker only when the observed syndrome still appears
consistent.

\subsubsection*{C2: balanced score proposal policy}

The second policy is the main score-based proposal rule.  For weights
$\lambda^{(2)}=(\lambda_s,\lambda_t,\lambda_i,\lambda_o)$ define
\[
S_X^{(2)}
=
\lambda_s s_X
-\lambda_t d_X
-\lambda_i(1-a_X)
+\lambda_o u_X.
\]
The configuration used in the current seed-10 experiments is
\[
\lambda^{(2)}=(1.0,\,1.0,\,2.0,\,0.5).
\]
The rule is
\[
C_2
=
\begin{cases}
B, & S_B^{(2)}>S_A^{(2)},\\
A, & S_B^{(2)}<S_A^{(2)}.
\end{cases}
\]
When admissibility differs and the two scores are exactly tied, the implementation
breaks the tie toward the admissible candidate; otherwise exact ties preserve $A$.
This small tie convention is included only for determinism.  The substantive point is
that $C_2$ is the policy that may actively trade syndrome support against relational
objective quality, trajectory proximity, and admissibility penalty.

\subsubsection*{C3: hard-admissibility ablation}

The original $C_3$ should now be interpreted as an ablation of the hypothesis that a
hard admissibility filter defines a distinct safety-first operating zone.  Its rule is
\[
C_3
=
\begin{cases}
A, & a_A=1,\ a_B=0,\\
B, & a_A=0,\ a_B=1,\\
\arg\max_{X\in\{A,B\}} S_X^{(3)}, & a_A=a_B,
\end{cases}
\]
with ties in the last line resolved toward $A$.  In the equal-admissibility case the
implementation uses
\[
S_X^{(3)}
=
0.8\,s_X
-1.2\,d_X
+0.3\,u_X.
\]
The configured inadmissibility weight for the $C_3$ score template is $5.0$, but the
hard admissibility branch removes the penalty term once both candidates have the same
admissibility status.  This is precisely why $C_3$ is not the final redesign: it tests
a Boolean version of safety, whereas the observed failure mode is a downside-risk
question inside imperfect syndrome evidence.

\subsubsection*{C3R: robust veto layer on C2}

The final policy used by the current code is $C_3^{\mathrm{R}}$, not $C_3$.  It treats
$C_2$ as a proposal policy and applies a one-sided veto to proposed switches to $B$.
Let
\[
\Delta S:=S_B^{(2)}-S_A^{(2)},
\qquad
\Delta v:=v_A-v_B.
\]
Imperfect syndrome information is summarized by
\[
u_{\mathrm{raw}}
=
(1-r_{\mathrm{obs}})
+p_{\mathrm{noise}}
+a_{\mathrm{amb}}
+p_{\mathrm{meas}}
+p_{\mathrm{reset}}
+\ell_{\mathrm{syn}},
\]
\[
u_{\mathrm{syn}}
=
\mathrm{clip}_{[0,1]}(u_{\mathrm{raw}}),
\]
where $r_{\mathrm{obs}}$ is the syndrome observation ratio, $p_{\mathrm{noise}}$ is
the syndrome-bit corruption probability, $a_{\mathrm{amb}}$ is the ambiguity level,
$p_{\mathrm{meas}}$ and $p_{\mathrm{reset}}$ are measurement/reset error probabilities,
and $\ell_{\mathrm{syn}}$ is measured syndrome information loss.  The raw value is
kept as a diagnostic; the clipped value is used by the gate.

The robust gate is
\[
C_3^{\mathrm{R}}=B
\quad\Longleftrightarrow\quad
\Bigl[C_2=B\Bigr]
\wedge
\Bigl[\Delta S\ge\delta_S\Bigr]
\wedge
\Bigl[v_A\ge\delta_A\Bigr]
\wedge
\Bigl[a_B=1,\ v_B\le\delta_B\Bigr]
\wedge
\Bigl[u_{\mathrm{syn}}\le u_{\max}\Bigr].
\]
Otherwise $C_3^{\mathrm{R}}=A$.  The current configuration is
\[
\delta_S=0,\qquad
\delta_A=0,\qquad
\delta_B=0,\qquad
u_{\max}=0.99.
\]
With these defaults the $v_A$ condition is non-restrictive because $v_A\ge0$ by
definition; it is retained as a tuning handle for stricter future gates.  The two
active safety requirements are that $B$ itself be admissible and that the clipped
syndrome uncertainty not be saturated.  Thus $u_{\mathrm{syn}}=1$ is treated as too
degraded to trust a relational switch.

Equivalently, define
\[
G_R
:=
\mathbf{1}[C_2=B]\,
\mathbf{1}[\Delta S\ge\delta_S]\,
\mathbf{1}[v_A\ge\delta_A]\,
\mathbf{1}[a_B=1,\ v_B\le\delta_B]\,
\mathbf{1}[u_{\mathrm{syn}}\le u_{\max}].
\]
Then
\[
C_3^{\mathrm{R}}
=
\begin{cases}
B, & G_R=1,\\
A, & G_R=0.
\end{cases}
\]
This gives the immediate invariants
\[
C_3^{\mathrm{R}}=B \Longrightarrow C_2=B,
\qquad
\chi_B(C_3^{\mathrm{R}})\le \chi_B(C_2),
\]
where $\chi_B(C)$ is the indicator that policy $C$ selects $B$.  The robust layer can
only block $C_2$ switches; it never creates additional switches that $C_2$ rejected.

\paragraph{Why $C_3^{\mathrm{R}}$ replaces $C_3$.}
The negative clean-map result for $C_3$ is not the main evidence for the redesign.
It is only an ablation result showing that Boolean admissibility alone does not create
a useful new policy region once inadmissibility has already been priced into $C_2$.
The actual redesign is the move from a hard admissibility selector to a one-sided
downside-risk veto on $C_2$.  This matters because the relevant failure mode is not
``choose the admissible candidate''; it is ``avoid trusting a relational switch when
the available syndrome evidence is too degraded to justify the downside risk.''  A
structurally admissible $B$ can still lower logical fidelity, and an observed syndrome
that looks consistent can hide a true failure boundary under partial or corrupted
observation.  $C_3^{\mathrm{R}}$ is therefore the paper-facing policy, while $C_3$ is
kept as a diagnostic ablation.

\subsection{Evaluation Metrics for Robust Switching}

The clean-syndrome regime map is useful for identifying whether a policy has an
independent operating zone.  It is not sufficient for judging $C_3^{\mathrm{R}}$,
because $C_3^{\mathrm{R}}$ is primarily an imperfect-information safety layer.  We
therefore separate policy quality, downside-risk, and switch-intervention metrics.

\subsubsection*{Policy quality metrics}

For each policy $C\in\{C_1,C_2,C_3,C_3^{\mathrm{R}}\}$:
\[
\overline{\Delta F}_C,
\qquad
P(\Delta F_C\ge 0),
\qquad
P(F_C\ge F_{\mathrm{success}}),
\qquad
P(C=B)
\]
measure mean fidelity gain, non-worsening rate, logical-success rate, and relational
candidate selection rate.  These are performance metrics, not by themselves robust
safety metrics: a policy may improve the mean while worsening the lower tail.

\subsubsection*{False-safe and tail-risk metrics}

The structural false-safe rate measures whether a policy chooses a structurally unsafe
candidate while the observed syndrome appears consistent:
\[
\mathrm{FS}_{\mathrm{struct}}(C)
=
P\!\left[
\mathrm{observed\ syndrome\ consistent}
\wedge
C\ \mathrm{inadmissible}
\right].
\]
This metric is useful, but it misses a failure mode that matters directly for recovery:
a switch can be structurally admissible but still lower logical fidelity.  We therefore
also use the fidelity false-safe rate
\[
\mathrm{FS}_{F}(C)
=
P\!\left[
\Delta F_C < -\varepsilon_F
\right],
\]
with $\varepsilon_F$ equal to the experiment's fidelity margin.

For magnitude-sensitive downside risk, define the empirical loss excess
\[
R_-(C)
:=
\mathbb{E}\!\left[
\max\{0,-\Delta F_C-\varepsilon_F\}
\right].
\]
We also track lower-tail statistics of the gain distribution,
\[
q_{.05}(C)
:=
\inf\{x:\widehat{P}(\Delta F_C\le x)\ge 0.05\},
\]
\[
\mathrm{CVaR}_{.05}(C)
:=
\widehat{\mathbb{E}}\!\left[
\Delta F_C\mid \Delta F_C\le q_{.05}(C)
\right].
\]
The intended conservative claim is therefore not that $C_3^{\mathrm{R}}$ uniformly
dominates $C_2$.  The claim is that the gate is explicitly parameterized to reduce
the frequency and magnitude of downside switches when syndrome information becomes
unreliable, and that this reduction should be visible in $\mathrm{FS}_F$, $q_{.05}$,
$\mathrm{CVaR}_{.05}$, and the intervention diagnostics below.

\subsubsection*{C3-R switch-intervention metrics}

Since $C_3^{\mathrm{R}}$ is only active when $C_2$ proposes $B$, the primitive block
indicator is
\[
B_R
:=
\mathbf{1}[C_2=B]\,
\mathbf{1}[C_3^{\mathrm{R}}=A].
\]
Two candidate-level ground-truth comparisons are then recorded:
\[
H_2
:=
\mathbf{1}[C_2=B]\,
\mathbf{1}[\Delta F_A>\Delta F_B+\varepsilon_F],
\]
\[
G_2
:=
\mathbf{1}[C_2=B]\,
\mathbf{1}[\Delta F_B>\Delta F_A+\varepsilon_F].
\]
Here $H_2$ means that the proposed $C_2$ switch to $B$ was harmful relative to staying
with $A$, while $G_2$ means that the proposed switch was beneficial.  The central
conditional rates are
\[
P(B_R=1\mid C_2=B),
\qquad
P(H_2=1\mid B_R=1),
\qquad
P(B_R=1\mid H_2=1),
\]
and
\[
P(B_R=1\mid G_2=1),
\qquad
P(B_R=0\mid G_2=1).
\]
We call these the block rate among $C_2$ switches, harmful-block precision,
harmful-switch recall, beneficial-switch block rate, and beneficial-switch retention.
This makes explicit that $C_3^{\mathrm{R}}$ is not a universal harmful-switch oracle:
it is an uncertainty-triggered safety gate that trades off prevented losses against
missed beneficial switches.

The corresponding magnitude decomposition is
\[
L_{\mathrm{prev}}
:=
\sum_i B_R^{(i)}
\max\{0,\Delta F_A^{(i)}-\Delta F_B^{(i)}\},
\]
\[
G_{\mathrm{miss}}
:=
\sum_i B_R^{(i)}
\max\{0,\Delta F_B^{(i)}-\Delta F_A^{(i)}\},
\]
\[
G_{\mathrm{net}}
:=
L_{\mathrm{prev}}-G_{\mathrm{miss}}
=
\sum_i\left(
\Delta F_{C_3^{\mathrm{R}}}^{(i)}
-
\Delta F_{C_2}^{(i)}
\right).
\]
Thus $G_{\mathrm{net}}>0$ says only that, on the sampled grid, blocked harmful-switch
losses outweighed missed beneficial-switch gains.  It is a sample-level intervention
diagnostic, not a theorem of dominance.

\subsubsection*{Oracle regret and failure-boundary metrics}

For each row define the two-candidate oracle gain
\[
\Delta F^\star
:=
\max\{\Delta F_A,\Delta F_B\}.
\]
The oracle regret of policy $C$ is
\[
\mathrm{Reg}(C)
:=
\Delta F^\star-\Delta F_C.
\]
Mean, $95$th-percentile, and upper-tail regret summarize how far the policy falls
from the best available candidate on the same row.  This is a stricter diagnostic
than comparing only $C_2$ and $C_3^{\mathrm{R}}$.

Finally, we distinguish observed and true failure boundaries:
\[
\mathrm{FB}_{\mathrm{obs}}(C)
:=
P\!\left[
\mathrm{observed\ syndrome\ consistent}
\wedge
\neg\,\mathrm{logical\ success}_C
\right],
\]
\[
\mathrm{FB}_{\mathrm{true}}(C)
:=
P\!\left[
\mathrm{true\ syndrome\ consistent}
\wedge
\neg\,\mathrm{logical\ success}_C
\right].
\]
The observed boundary is what the policy can directly see.  The true boundary is an
offline diagnostic used to check whether observed consistency is hiding a logical
failure under partial, noisy, or ambiguous syndrome information.

\subsection{Experimental Stress Protocols}

The hybrid layer is evaluated under five increasingly degraded syndrome regimes.
\begin{itemize}
\item \textbf{Clean regime map.}  Complete, uncorrupted syndrome observations.  This
tests whether a policy has an independent operating zone when syndrome information is
reliable.
\item \textbf{Partial syndrome.}  Only a fraction $r_{\mathrm{obs}}$ of syndrome
information is exposed to the policy layer.  This stresses uncertainty from missing
syndrome information.
\item \textbf{Noisy syndrome.}  Syndrome bits are corrupted with probability
$p_{\mathrm{noise}}$.  This stresses random observational corruption while keeping
the observation ratio fixed.
\item \textbf{Partial+noisy syndrome.}  Missing and corrupted syndrome information
are combined.  This is the most direct stress test for over-trusting relational
switches under degraded evidence.
\item \textbf{Ambiguity+measurement/reset syndrome.}  Syndrome ambiguity and physical
measurement/reset error probabilities are swept jointly.  This isolates whether the
robust gate reacts to ambiguity-driven uncertainty rather than only to missing bits.
\end{itemize}

The clean regime map and the imperfect-syndrome maps answer different questions.
The clean map asks whether a policy defines a distinct operating zone.  The imperfect
maps ask whether a policy avoids downside risk when the evidence for switching has
become less trustworthy.

\subsection{Empirical Operating Zones of $C_1$, $C_2$, $C_3$, and $C_3^{\mathrm{R}}$}

In the clean-syndrome seed-10 regime map, the full encoded grid is aggregated over
$240$ regime cells after averaging over seeds.  The preferred-policy counts are
\[
\#(C_1)=197,
\qquad
\#(C_2)=43,
\qquad
\#(C_3)=0,
\qquad
\#(C_3^{\mathrm{R}})=0.
\]
Thus $C_1$ remains the dominant default policy, $C_2$ occupies a nontrivial local
operating zone, and neither $C_3$ nor $C_3^{\mathrm{R}}$ forms an independent region
on the clean-syndrome grid.  This is expected for $C_3^{\mathrm{R}}$: when syndrome
observation is complete and uncorrupted, the raw uncertainty is zero, the uncertainty
gate is open, and the robust policy collapses to the same choices as $C_2$.

The $C_2$-preferred cells form a structured cluster rather than a scattered set of
exceptions.  All $43$ $C_2$-preferred clean-regime cells occur under the phaseflip
code, and they are concentrated in phase-sensitive stochastic sectors:
\[
16\ \mathrm{dephasing},\qquad
10\ \mathrm{phaseflip},\qquad
9\ \mathrm{mixed\ Pauli},\qquad
6\ \mathrm{depolarizing},\qquad
2\ \mathrm{bitflip}.
\]
In these $C_2$-preferred cells the mean logical-fidelity gains satisfy
\[
\overline{\Delta F}_{C_1}\approx 0.0135,
\qquad
\overline{\Delta F}_{C_2}\approx 0.0250,
\]
while the mean rates of selecting the relational candidate are
\[
\overline{\chi_B(C_1)}\approx 0.9349,
\qquad
\overline{\chi_B(C_2)}\approx 0.4000.
\]
The fidelity false-safe rates in the same cells are
\[
\overline{\mathrm{FS}_{F}(C_1)}\approx 0.2674,
\qquad
\overline{\mathrm{FS}_{F}(C_2)}\approx 0.1093.
\]
Hence the improvement of $C_2$ comes from switching more selectively, not from always
trusting the relational candidate.  The rule-based $C_1$ policy tends to over-select
$B$ near phase-sensitive boundaries where the syndrome baseline should still be
preserved.

\subsection{Final Seed-10 C3-R Results}

The final seed-10 rerun was evaluated under the clipped uncertainty convention and
$u_{\max}=0.99$.  The resulting policy-level outcomes are:

\begin{center}
\scriptsize
\begin{tabular}{l|r|cc|cc|cc|cc}
Regime
& $n$
& $\overline{\Delta F}_{C_2}$
& $\overline{\Delta F}_{C_3^{\mathrm{R}}}$
& $P_{\mathrm{succ}}(C_2)$
& $P_{\mathrm{succ}}(C_3^{\mathrm{R}})$
& $\mathrm{FS}_F(C_2)$
& $\mathrm{FS}_F(C_3^{\mathrm{R}})$
& $q_{.05}(C_2)$
& $q_{.05}(C_3^{\mathrm{R}})$
\\
\hline
clean & 2400 & 0.1273 & 0.1273 & 0.6358 & 0.6358 & 0.0238 & 0.0238 & 0.0000 & 0.0000\\
partial & 2880 & 0.1169 & 0.1189 & 0.4552 & 0.5976 & 0.0694 & 0.0389 & -0.0152 & -0.0006\\
noisy & 3600 & 0.1169 & 0.1169 & 0.4556 & 0.4556 & 0.0678 & 0.0678 & -0.0104 & -0.0104\\
partial+noisy & 2160 & 0.1170 & 0.1212 & 0.4551 & 0.7241 & 0.0685 & 0.0134 & -0.0104 & 0.0000\\
ambiguity+meas. & 3600 & 0.1600 & 0.1604 & 0.5208 & 0.5978 & 0.0000 & 0.0000 & 0.0192 & 0.0238
\end{tabular}
\end{center}

The lower-tail and oracle diagnostics give the same qualitative picture:

\begin{center}
\scriptsize
\begin{tabular}{l|cc|cc|cc}
Regime
& $\mathrm{CVaR}_{.05}(C_2)$
& $\mathrm{CVaR}_{.05}(C_3^{\mathrm{R}})$
& $\overline{\mathrm{Reg}}(C_2)$
& $\overline{\mathrm{Reg}}(C_3^{\mathrm{R}})$
& $\mathrm{FB}_{\mathrm{true}}(C_2)$
& $\mathrm{FB}_{\mathrm{true}}(C_3^{\mathrm{R}})$
\\
\hline
clean & -0.0033 & -0.0033 & 0.0132 & 0.0132 & 0.2375 & 0.2375\\
partial & -0.0208 & -0.0147 & 0.0271 & 0.0251 & 0.3611 & 0.2653\\
noisy & -0.0207 & -0.0207 & 0.0270 & 0.0270 & 0.3589 & 0.3589\\
partial+noisy & -0.0206 & -0.0034 & 0.0271 & 0.0228 & 0.3616 & 0.1819\\
ambiguity+meas. & 0.0078 & 0.0084 & 0.0336 & 0.0332 & 0.2136 & 0.1847
\end{tabular}
\end{center}

The switch-intervention table is the most direct evidence for the redesigned policy:

\begin{center}
\scriptsize
\begin{tabular}{l|r|r|ccccc|r}
Regime
& $C_2=B$
& blocked
& $P(B_R\mid C_2=B)$
& precision
& recall
& ben. block
& ben. retain
& $G_{\mathrm{net}}$
\\
\hline
clean & 1529 & 0 & 0.0000 & -- & -- & -- & -- & 0.0000\\
partial & 1094 & 476 & 0.4351 & 0.8697 & 0.4461 & 0.3750 & 0.6250 & 5.6616\\
noisy & 1370 & 0 & 0.0000 & -- & -- & -- & -- & 0.0000\\
partial+noisy & 821 & 670 & 0.8161 & 0.8791 & 0.8475 & 0.6739 & 0.3261 & 9.1212\\
ambiguity+meas. & 895 & 357 & 0.3989 & 0.7759 & 0.3929 & 0.4211 & 0.5789 & 1.5749
\end{tabular}
\end{center}

\paragraph{Interpretation.}
The clean map confirms that $C_3^{\mathrm{R}}$ has no separate clean-syndrome operating
zone; it reduces to $C_2$ because the uncertainty evidence does not ask for a veto.
The noisy-only map also leaves $C_3^{\mathrm{R}}$ inactive.  This is not a failure of
the implementation: under this stress protocol the raw uncertainty mean is only
$0.038$, far below the saturation threshold, so the robust gate has no reason to block
$C_2$ switches.

The partial and partial+noisy maps are the main positive evidence.  In the partial
map, $C_3^{\mathrm{R}}$ blocks $43.51\%$ of $C_2$ switches, with harmful-block
precision $86.97\%$ and net intervention gain $5.6616$.  In the partial+noisy map,
the effect is stronger: $C_3^{\mathrm{R}}$ blocks $81.61\%$ of $C_2$ switches,
recovers the lower tail from $q_{.05}=-0.0104$ to $0.0000$, improves
$\mathrm{CVaR}_{.05}$ from $-0.0206$ to $-0.0034$, and reduces the true failure
boundary from $0.3616$ to $0.1819$.

The ambiguity+measurement map shows why the old false-safe-only evaluation was too
narrow.  Both $C_2$ and $C_3^{\mathrm{R}}$ have $\mathrm{FS}_F=0$, but the robust gate
still slightly improves mean gain, logical success, lower-tail gain, oracle regret,
and true failure-boundary rate.  These improvements are modest rather than decisive,
so this regime should be described as supporting evidence for the safety-gate
interpretation, not as a separate dominance claim.

The main cost is also visible.  $C_3^{\mathrm{R}}$ sometimes blocks beneficial $C_2$
switches: the beneficial-switch block rate is $37.50\%$ in the partial map,
$67.39\%$ in the partial+noisy map, and $42.11\%$ in the ambiguity+measurement map.
Thus the conservative paper-facing statement should be that $C_3^{\mathrm{R}}$ is a
downside-risk policy whose sampled net intervention gain is positive in the degraded
syndrome regimes tested here.  It should not yet be stated as an optimal or uniformly
dominant switching policy.

\paragraph{Taxonomy after the redesign.}
The present interpretation is:
\begin{itemize}
\item $C_1$: default rule-based hybrid on most clean regimes,
\item $C_2$: balanced phase-sensitive correction that suppresses over-switching,
\item $C_3$: hard-admissibility probe with no independent observed operating zone,
\item $C_3^{\mathrm{R}}$: no independent clean-syndrome zone, but a dedicated
imperfect-information safety layer that vetoes selected $C_2$ switches under degraded
syndrome evidence.
\end{itemize}
The role of $C_3^{\mathrm{R}}$ is therefore not to outperform $C_2$ in every clean or
low-uncertainty setting.  Its role is to reduce downside-risk exposure when the
evidence for switching becomes partial, ambiguous, or jointly corrupted.

\section{Clock Geometry and Relational Consistency}

\subsection{Clock Orbit and Local Distinguishability}

Let the clock states be generated by
\begin{equation}
|\tau\rangle_C = e^{-i\hat{H}_C \tau} |0\rangle_C,
\end{equation}
where $\hat{H}_C$ is the clock Hamiltonian.

We define the local distinguishability between adjacent clock labels as
\begin{equation}
d_{\mathrm{clock}}^{\mathrm{geom}}(\tau,\tau+\delta\tau)
:=
1 - \left| {}_C\langle \tau \mid \tau+\delta\tau \rangle_C \right|^2.
\end{equation}

Equivalently, using the Bures distance,
\begin{equation}
D_{\mathrm{B}}(\tau,\tau+\delta\tau)
=
\sqrt{2\left(1 - \sqrt{F_\tau}\right)},
\end{equation}
where
\begin{equation}
F_\tau
=
F\!\left(
|\tau\rangle\langle\tau|_C,
|\tau+\delta\tau\rangle\langle\tau+\delta\tau|_C
\right)
=
\left| {}_C\langle \tau \mid \tau+\delta\tau \rangle_C \right|^2.
\end{equation}

For sufficiently small $\delta\tau$, we obtain the expansion
\begin{equation}
d_{\mathrm{clock}}^{\mathrm{geom}}(\tau,\tau+\delta\tau)
=
(\Delta H_C)^2_\tau\,(\delta\tau)^2
+ O\!\left((\delta\tau)^4\right),
\end{equation}
where
\begin{equation}
(\Delta H_C)^2_\tau
=
\langle \tau | \hat{H}_C^2 | \tau \rangle
-
\langle \tau | \hat{H}_C | \tau \rangle^2.
\end{equation}

\paragraph{Interpretation}
The quantity $d_{\mathrm{clock}}^{\mathrm{geom}}$ measures the intrinsic distinguishability of adjacent clock labels and induces a local geometric structure on the clock orbit. It does not depend on external time and serves as an intrinsic resolution measure of the relational clock.

\subsection{Relational Response Between Clock Slices}

Let $\rho_S(\tau)$ be the clock-conditioned system state.
We define the relational response between adjacent slices as
\begin{equation}
d_{\mathrm{resp}}(\tau,\tau+\delta\tau)
:=
d\!\left(
\rho_S(\tau), \rho_S(\tau+\delta\tau)
\right),
\end{equation}
In the main formulation, we instantiate $d_{\mathrm{resp}}$ as 
the spectral response
\[
d_{\mathrm{resp}}(\tau, \tau+\delta\tau) 
:= \|\Lambda(\tau+\delta\tau) - \Lambda(\tau)\|_2,
\]
thereby aligning the clock-consistency penalty with the geometry 
of the spectral trajectory space. The fidelity-based alternative 
is discussed in the appendix as a robustness check.

\subsection{Clock Consistency Measure}

We define the clock consistency ratio as
\begin{equation}
\mathcal{C}_{\mathrm{clock}}(\tau)
=
\frac{
d_{\mathrm{resp}}(\tau,\tau+\delta\tau)
}{
d_{\mathrm{clock}}^{\mathrm{geom}}(\tau,\tau+\delta\tau)
+ \epsilon
},
\end{equation}
where $\epsilon > 0$ is a small regularization constant.

\paragraph{Interpretation}
This ratio measures whether the variation of the system state is consistent with the intrinsic resolution of the clock.

\begin{itemize}
\item $\mathcal{C}_{\mathrm{clock}}(\tau) \ll 1$: stable relational encoding
\item $\mathcal{C}_{\mathrm{clock}}(\tau) \approx 1$: consistent evolution
\item $\mathcal{C}_{\mathrm{clock}}(\tau) \gg 1$: anomalous or inconsistent slicing
\end{itemize}

\subsection{Trajectory-Level Clock Consistency}

We extend the definition to the full trajectory:
\begin{equation}
\Phi_{\mathrm{clock}}(\Gamma)
=
\int_{\mathcal{T}}
\mathcal{C}_{\mathrm{clock}}(\tau)
\, d\mu(\tau),
\end{equation}
or in discrete form,
\begin{equation}
\Phi_{\mathrm{clock}}(\Gamma)
=
\sum_{r=1}^{m-1}
\frac{
d\!\left(\rho_S(\tau_r),\rho_S(\tau_{r+1})\right)
}{
d_{\mathrm{clock}}^{\mathrm{geom}}(\tau_r,\tau_{r+1}) + \epsilon
}.
\end{equation}

A trajectory is considered clock-consistent if
\begin{equation}
\Phi_{\mathrm{clock}}(\Gamma) \le \eta.
\end{equation}

\begin{remark}
The admissibility threshold $\eta$ is calibrated statistically 
from the background distribution of $\Phi_{\mathrm{clock}}$ 
evaluated on unperturbed reference trajectories, whereas the 
slice spacing $\delta\tau$ is constrained by the clock resolution 
bound of Lemma~4.1. These two quantities serve distinct roles 
and are not directly related.
\end{remark}

The abstract clock penalty $\Phi_{\mathrm{clock}}$ introduced in 
Eq.~(89) is concretely instantiated in Sec.~11 through the 
clock-geometry-based consistency ratio in Eqs.~(99)--(101).

\end{document}
